\begin{table*}[t]
    \centering
    \small
    \sisetup{table-number-alignment=center}
    \begin{tabularx}{\textwidth}{|l|l|l|l|X|}
      \hline
      \textbf{Signal} & \textbf{Range} & \textbf{Components} & \textbf{Elements} & \textbf{Rationale for inclusion} \\
      \hline
      Cap sensing (pressure) & \(0 \dots 800\) & 18 electrodes & Ele\_0--Ele\_17 & High-resolution plantar pressure map. Captures stance phase, load distribution, and contact pattern differences across activities (e.g., forefoot vs.\ heel loading in jumping vs.\ walking). \\
      \hline
      Accelerometer (linear) & \(-1 \dots 1\) & \(x, y, z\) & Ele\_18--Ele\_20 & Foot-frame linear acceleration. Encodes impact magnitude, swing dynamics, and vertical motion, which are highly discriminative for gait phase transitions and dynamic activities such as jumping. \\
      \hline
      Quaternion (orientation) & \(-1 \dots 1\) & \(x, y, z, v\) & Ele\_30--Ele\_33 & 3D orientation of the insole, fused from inertial and magnetic sensors. Provides a numerically stable, gimbal-lock-free representation of foot pose over time, enabling the model to learn differences in foot rotation patterns between gait types. \\
      \hline
    \end{tabularx}
    \caption{Subset of insole signals retained for model input. We keep per-timestep plantar pressure, linear acceleration, and fused orientation (quaternions) for each foot, as these jointly capture contact mechanics, dynamic motion, and foot pose. Redundant timestamps, raw gyroscope/magnetometer channels, and battery state are excluded from the learning pipeline.}
    \label{tab:insole-features-selected}
  \end{table*}

